\chaptersection{3}{Failure Mode and Effects Analysis}

La Failure Mode and Effects Analysis (FMEA) è una metodologia utilizzata
per identificare le possibili modalità di guasto di un sistema e valutarne
le conseguenze. L'analisi considera gli asset selezionati tra quelli
individuati nell'Asset Model, concentrandosi sugli elementi maggiormente
significativi per il funzionamento della pipeline di DermaTriage. Per
ciascun asset vengono esaminate le funzioni svolte, i relativi prerequisiti
e i possibili failure mode, individuandone le cause e gli effetti che
possono propagarsi agli altri componenti del sistema.

\subsection{A01 -- Immagine dermatologica}

L'asset A01 rappresenta l'immagine della lesione dermatologica acquisita
per il caso clinico e utilizzata come principale informazione visiva
della pipeline. L'immagine può essere recuperata nell'ambito del consulto
gestito tramite B4 oppure fornita direttamente attraverso l'endpoint
\texttt{/analyze}.

All'interno della pipeline, l'immagine viene utilizzata nello Stage 1
da EfficientNet-B4 per la classificazione iniziale dell'urgenza e nello
Stage 2 da Qwen2-VL per la generazione della descrizione clinica della
lesione. La correttezza e l'adeguatezza dell'immagine costituiscono quindi
un prerequisito per le successive elaborazioni basate sull'informazione
visiva.


\subsubsection{Funzioni e prerequisiti}

Le principali funzioni dell'asset e i prerequisiti necessari al loro
svolgimento sono riportati nella Tabella~\ref{tab:a01-functions}.

\begin{table}[H]
    \centering
    \small
    \setstretch{1}
    \setlength{\tabcolsep}{4pt}
    \renewcommand{\arraystretch}{1.08}

    \begin{tabularx}{\textwidth}{
        >{\raggedright\arraybackslash}p{1.1cm}
        >{\raggedright\arraybackslash}p{3.8cm}
        X
    }
        \toprule

        \textbf{ID} &
        \textbf{Funzione} &
        \textbf{Prerequisiti} \\

        \midrule

        F01.1 &
        Rappresentazione visiva della lesione &
        Acquisizione di un'immagine che rappresenti in modo adeguato
        la lesione dermatologica oggetto del consulto. \\

        \midrule

        F01.2 &
        Associazione dell'immagine al caso clinico &
        Corretta corrispondenza tra l'immagine acquisita e il consulto
        o la richiesta di analisi a cui essa appartiene. \\

        \midrule

        F01.3 &
        Fornitura dell'input visivo alla pipeline di inferenza &
        Disponibilità dell'immagine in un formato utilizzabile e corretta
        trasmissione dell'input a EfficientNet-B4 nello Stage 1 e a Qwen2-VL
        nello Stage 2. \\

        \bottomrule
    \end{tabularx}

    \caption{Funzioni e prerequisiti dell'asset A01 -- Immagine dermatologica.}
    \label{tab:a01-functions}
\end{table}


\subsubsection{Failure mode}

Sulla base delle funzioni individuate vengono considerate le principali
modalità con cui l'asset A01 può non svolgere correttamente il proprio
compito, ovvero i failure mode.

\begin{table}[H]
    \centering
    \footnotesize
    \setstretch{1}
    \setlength{\tabcolsep}{3pt}
    \renewcommand{\arraystretch}{1.1}

    \begin{tabularx}{\textwidth}{
        >{\raggedright\arraybackslash}p{1.1cm}
        >{\raggedright\arraybackslash}p{3cm}
        >{\raggedright\arraybackslash}p{5cm}
        X
    }
        \toprule

        \textbf{ID} &
        \textbf{Failure mode} &
        \textbf{Possibili cause} &
        \textbf{Effetti} \\

        \midrule

        FM01.1 &
        Immagine di qualità insufficiente &
        Sfocatura, illuminazione inadeguata, risoluzione insufficiente,
        presenza di artefatti oppure condizioni di acquisizione che
        rendono poco distinguibili le caratteristiche della lesione. &
        L'immagine non rappresenta con sufficiente chiarezza le
        caratteristiche visive della lesione, riducendo l'affidabilità
        delle elaborazioni effettuate da EfficientNet-B4 e Qwen2-VL. \\

        \midrule

        FM01.2 &
        Lesione rappresentata in modo incompleto o non adeguato &
        Inquadratura parziale, porzione rilevante della lesione esclusa
        dall'immagine oppure acquisizione di una regione cutanea che non
        consente di osservare adeguatamente la lesione di interesse. &
        I modelli ricevono una rappresentazione visiva incompleta del caso
        e possono quindi produrre una classificazione dell'urgenza o una
        descrizione clinica non rappresentative della lesione effettiva. \\

        \midrule

        FM01.3 &
        Immagine associata al caso clinico errato &
        Errata associazione tra immagine e consulto, selezione del
        documento non corrispondente al caso oppure errore nella gestione
        del riferimento utilizzato per recuperare l'immagine. &
        La pipeline elabora un'immagine appartenente a un caso differente
        rispetto alle informazioni cliniche associate alla richiesta.
        Le inferenze prodotte risultano quindi riferite a dati non
        coerenti tra loro e possono determinare un risultato di triage
        non pertinente al paziente analizzato. \\

        \midrule

        FM01.4 &
        Immagine non disponibile o non utilizzabile dalla pipeline &
        Formato non supportato, file danneggiato, errore nella lettura del file,
        mancato recupero dell'immagine dal servizio di origine oppure errore
        durante il trasferimento. &
        EfficientNet-B4 e Qwen2-VL non dispongono di un input visivo utilizzabile
        e non possono quindi completare normalmente le rispettive elaborazioni. \\

        \bottomrule
    \end{tabularx}

    \caption{Failure mode individuati per l'asset A01 -- Immagine dermatologica.}
    \label{tab:a01-failure-modes}
\end{table}
\subsection{A09 -- Dataset clinico RAG}

L'asset A09 rappresenta il dataset clinico utilizzato dal sistema RAG,
memorizzato nel file \texttt{RAG\_dataset.csv}. Il dataset contiene casi
clinici dermatologici caratterizzati da informazioni demografiche,
localizzazione della lesione, descrizione clinica, diagnosi e livello
di urgenza.

Durante l'inferenza, le descrizioni cliniche contenute nel dataset
costituiscono la base informativa utilizzata nello Stage 3 per il recupero
dei casi clinici più simili alla descrizione prodotta da Qwen2-VL.
Il dataset rappresenta inoltre la base da cui vengono derivati i dati
utilizzati per l'addestramento del classificatore EfficientNet-B4.


\subsubsection{Funzioni e prerequisiti}

Le principali funzioni dell'asset e i prerequisiti necessari al loro
svolgimento sono riportati nella Tabella~\ref{tab:a09-functions}.

\begin{table}[H]
    \centering
    \small
    \setstretch{1}
    \setlength{\tabcolsep}{4pt}
    \renewcommand{\arraystretch}{1.08}

    \begin{tabularx}{\textwidth}{
        >{\raggedright\arraybackslash}p{1.1cm}
        >{\raggedright\arraybackslash}p{3.8cm}
        X
    }
        \toprule

        \textbf{ID} &
        \textbf{Funzione} &
        \textbf{Prerequisiti} \\

        \midrule

        F09.1 &
        Rappresentazione dei casi clinici dermatologici &
        Presenza di record clinici correttamente strutturati e contenenti
        informazioni coerenti relative a descrizione della lesione,
        diagnosi e livello di urgenza. \\

        \midrule

        F09.2 &
        Fornitura della base informativa per il retrieval RAG &
        Disponibilità delle descrizioni cliniche necessarie
        all'indicizzazione e corretta associazione tra descrizione,
        diagnosi, urgenza e relativo caso clinico. \\

        \midrule

        F09.3 &
        Supporto alla costruzione dei dati per l'addestramento &
        Disponibilità di record e label utilizzabili per derivare
        i dataset impiegati nel processo di training del classificatore. \\

        \bottomrule
    \end{tabularx}

    \caption{Funzioni e prerequisiti dell'asset A09 -- Dataset clinico RAG.}
    \label{tab:a09-functions}
\end{table}


\subsubsection{Failure mode}

Sulla base delle funzioni individuate vengono considerate le principali
modalità con cui l'asset A09 può non svolgere correttamente il proprio
compito, ovvero i failure mode.

\begin{table}[H]
    \centering
    \footnotesize
    \setstretch{1}
    \setlength{\tabcolsep}{3pt}
    \renewcommand{\arraystretch}{1.1}

    \begin{tabularx}{\textwidth}{
        >{\raggedright\arraybackslash}p{1.1cm}
        >{\raggedright\arraybackslash}p{3cm}
        >{\raggedright\arraybackslash}p{5cm}
        X
    }
        \toprule

        \textbf{ID} &
        \textbf{Failure mode} &
        \textbf{Possibili cause} &
        \textbf{Effetti} \\

        \midrule

        FM09.1 &
        Record clinici incompleti o incoerenti &
        Valori mancanti, informazioni cliniche non coerenti tra loro,
        errori nella compilazione dei record oppure associazione errata
        tra attributi appartenenti allo stesso caso. &
        Il dataset rappresenta in modo incompleto o non corretto i casi
        clinici e può fornire informazioni non affidabili sia al sistema
        RAG sia ai processi che utilizzano il dataset come base per la
        preparazione dei dati di addestramento. \\

        \midrule

        FM09.2 &
        Associazione errata tra descrizione clinica e diagnosi o urgenza &
        Label diagnostica o livello di urgenza non corrispondenti alla
        descrizione clinica del caso oppure errore nell'associazione dei
        campi appartenenti allo stesso record. &
        Un caso clinico può essere recuperato come riferimento insieme
        a una diagnosi o a un livello di urgenza non corretti. L'errore
        può propagarsi allo Stage 3 e influenzare il contesto fornito
        a BioMistral, oppure essere trasferito ai dati utilizzati
        nell'addestramento del classificatore. \\

        \midrule

        FM09.3 &
        Dataset non sufficientemente rappresentativo del dominio &
        Distribuzione sbilanciata dei casi, sottorappresentazione di
        determinate tipologie di lesione o caratteristiche della
        popolazione non adeguatamente coperte dal dataset. &
        Il sistema dispone di una base clinica meno rappresentativa
        rispetto ai casi incontrati in esercizio. Il retrieval può
        individuare riferimenti poco pertinenti per alcuni casi e i dati
        derivati dal dataset possono contribuire a una capacità di
        generalizzazione insufficiente del classificatore. \\

        \midrule

        FM09.4 &
        Dataset non disponibile o non correttamente accessibile &
        File mancante o danneggiato, percorso configurato in modo errato,
        errore durante la lettura oppure indisponibilità della risorsa
        nel momento in cui deve essere utilizzata. &
        Il sistema non dispone della base informativa necessaria per
        costruire o aggiornare il supporto al retrieval RAG e i processi
        che dipendono dal dataset non possono svolgersi secondo il normale
        funzionamento previsto. \\

        \bottomrule
    \end{tabularx}

    \caption{Failure mode individuati per l'asset A09 -- Dataset clinico RAG.}
    \label{tab:a09-failure-modes}
\end{table}
\subsection{A16 -- Modello EfficientNet-B4}

L'asset A16 riguarda il modello EfficientNet-B4 utilizzato per la
classificazione iniziale dell'urgenza dermatologica e i parametri appresi
durante l'addestramento, caricati dal checkpoint attivo. A partire
dall'immagine della lesione, il modello produce una classe di urgenza
tra HIGH, MEDIUM e LOW, insieme alla confidence della classificazione e
alle probabilità associate alle diverse classi.

Il risultato dello Stage 1 viene utilizzato dalle successive elaborazioni
della pipeline e contribuisce alla sintesi clinica prodotta da BioMistral
e alla successiva determinazione della priorità P1--P4.


\subsubsection{Funzioni e prerequisiti}

Le principali funzioni dell'asset e i prerequisiti necessari al loro
svolgimento sono riportati nella Tabella~\ref{tab:a16-functions}.

\begin{table}[H]
    \centering
    \small
    \setstretch{1}
    \setlength{\tabcolsep}{4pt}
    \renewcommand{\arraystretch}{1.08}

    \begin{tabularx}{\textwidth}{
        >{\raggedright\arraybackslash}p{1.1cm}
        >{\raggedright\arraybackslash}p{3.8cm}
        X
    }
        \toprule

        \textbf{ID} &
        \textbf{Funzione} &
        \textbf{Prerequisiti} \\

        \midrule

        F16.1 &
        Estrazione delle caratteristiche visive &
        Disponibilità di un'immagine dermatologica acquisita correttamente
        (A01) e caricamento del modello con i relativi parametri. \\

        \midrule

        F16.2 &
        Classificazione dell'urgenza in HIGH, MEDIUM o LOW &
        Corretta elaborazione dell'immagine e utilizzo dei parametri
        previsti per la versione attiva del classificatore. \\

        \midrule

        F16.3 &
        Produzione della confidence e delle probabilità associate &
        Completamento corretto dell'inferenza e disponibilità dell'output
        numerico prodotto dal classificatore. \\

        \midrule

        F16.4 &
        Utilizzo della versione prevista del modello &
        Disponibilità del checkpoint attivo e corretta configurazione
        dei riferimenti utilizzati per il caricamento del modello (A23). \\

        \bottomrule
    \end{tabularx}

    \caption{Funzioni e prerequisiti dell'asset A16 -- Modello EfficientNet-B4 e relativi pesi.}
    \label{tab:a16-functions}
\end{table}


\subsubsection{Failure mode}

Sulla base delle funzioni individuate vengono considerate le principali
modalità con cui l'asset A16 può non svolgere correttamente il proprio
compito, ovvero i failure mode.

\begin{table}[H]
    \centering
    \footnotesize
    \setstretch{1}
    \setlength{\tabcolsep}{3pt}
    \renewcommand{\arraystretch}{1.1}

    \begin{tabularx}{\textwidth}{
        >{\raggedright\arraybackslash}p{1.1cm}
        >{\raggedright\arraybackslash}p{3cm}
        >{\raggedright\arraybackslash}p{5cm}
        X
    }
        \toprule

        \textbf{ID} &
        \textbf{Failure mode} &
        \textbf{Possibili cause} &
        \textbf{Effetti} \\

        \midrule

        FM16.1 &
        Estrazione inadeguata delle caratteristiche visive &
        Qualità insufficiente dell'immagine, presenza di artefatti visivi oppure
        inquadratura non adeguata della lesione. &
        Le caratteristiche ricavate dall'immagine non rappresentano
        adeguatamente la lesione osservata, compromettendo l'informazione
        utilizzata dal classificatore per determinare la classe di urgenza. \\

        \midrule

        FM16.2 &
        Classificazione errata dell'urgenza &
        Limiti di generalizzazione del modello, distribuzione dei casi operativi
        differente da quella osservata durante l'addestramento oppure insufficiente
        capacità discriminativa tra le classi di urgenza. &
        Il modello assegna una classe HIGH, MEDIUM o LOW non coerente con
        l'effettiva urgenza del caso. L'errore può determinare una sovrastima
        o una sottostima dell'urgenza che si propaga agli stadi
        successivi, contribuendo così ad una valutazione finale del triage
        non appropriata. \\

        \midrule

        FM16.3 &
        Confidence o probabilità non coerenti con la classificazione &
        Calibrazione non adeguata delle probabilità prodotte dal modello
        oppure comportamento eccessivamente confidente su campioni
        incerti o poco rappresentativi. &
        La classe restituita è accompagnata da valori che non riflettono
        correttamente l'incertezza dell'inferenza. Gli stadi successivi
        ricevono quindi una valutazione non attendibile della confidenza
        associata alla classificazione. \\

        \midrule

        FM16.4 &
        Utilizzo di pesi o checkpoint non validi o non previsti &
        Corruzione o incompletezza del checkpoint, incompatibilità tra
        modello e parametri caricati, selezione di un checkpoint errato
        oppure mancato allineamento tra la versione validata e quella
        effettivamente utilizzata durante l'inferenza. &
        Il classificatore può operare con un comportamento diverso da
        quello previsto o validato, producendo classificazioni non
        affidabili. Nei casi in cui i parametri risultino incompatibili
        o non caricabili, l'inferenza può inoltre non essere completata. \\

        \midrule

        FM16.5 &
        Output di inferenza incompleto o non disponibile &
        Errore durante l'esecuzione del modello, incompatibilità tra input,
        modello e parametri caricati oppure anomalia software che interrompe
        l'elaborazione prima del suo completamento. &
        Lo Stage 1 non restituisce tutti gli elementi previsti, quali classe,
        confidence e probabilità, oppure non produce alcun risultato utilizzabile.
        Gli stadi dipendenti dall'output del classificatore non dispongono quindi
        delle informazioni necessarie per proseguire secondo il normale flusso
        della pipeline. \\

        \bottomrule
    \end{tabularx}

    \caption{Failure mode individuati per l'asset A16 -- Modello EfficientNet-B4 e relativi pesi.}
    \label{tab:a16-failure-modes}
\end{table}